% biology — том биологии, генетического кода, фолдинга и онтогенеза на кубооктаэдре ГЦК
\part{Биология, генетический код, фолдинг и онтогенез}

\chapter{Цифровая геометрия ДНК/РНК, замок Уотсона--Крика и штрихкод большой бороздки}
\label{ch:bio-dna-barcode}

\section{Вход из тома химии: изоэнергетический резонанс воды и соизмеримость пептидного шага}
\label{sec:bio-chem-input}

Биологический этаж не вводит виталистических правил или статистических подгонок поверх физики и химии
(\bcref{ch:alpha}, \bcref{ch:si-sm}): вся молекулярная биология опирается на три точных геометрических тождества
второй координационной сферы кубооктаэдра~\latGCKshort\ ($k=2$, $N_{\mathrm{shell}}(2)=10\cdot 2^2+2=42$ узла):
\begin{enumerate}[noitemsep]
\item \textbf{Резонанс кислорода и водорода ($Z=8\leftrightarrow Z=1$):}
  шесть квадратных граней канала $\square=(4,4,4)$ экранируют в атоме кислорода ровно $\sigma(\text{O})=|\square|=6$ единиц заряда,
  оставляя эффективный заряд $Z_{\mathrm{eff}}(\text{O})=8-6=2=n$ и делая первый потенциал ионизации кислорода
  равным боровской энергии связи водорода:
  \[
    IE_1(\text{O}) = \left(\frac{2}{2}\right)^{\!2} BE_H = BE_H = IE_1(\text{H}) \approx 13{,}61\ \mathrm{эВ}.
  \]
\item \textbf{Квант водородной связи $E_{\mathrm{HB}}$ и четырёхстороннее тождество водородного делителя:}
  на одну спиновую пару Паули второй координационной сферы ($\caKappaGeom^2 N_{\mathrm{shell}}(2)$ пар при $N_{\mathrm{shell}}(2)=(\geomNeighborTwelve-2)\cdot 2^2+2$)
  от ковалентной энергии $D_0(\text{H}_2)=BE_H/d$ приходится дискретный квант водородного замка, делитель которого удовлетворяет тождеству констант носителя:
  \[
    N_{\mathrm{HB}}
    := d\,\caKappaGeom^2 N_{\mathrm{shell}}(2)
    = \frac{1-\caKappaGeom^6}{\caKappaGeom^6}\,\lfloor\caBitBudget\rfloor
    = \vevHierarchySteps^2 - 1
    = 2^{\mathrm{frac\_bits}} - 1
    = z_{\min}^{-1} - 1,
  \]
  откуда без единого голого числа:
  \begin{equation}
  \label{eq:bio-hbond-quantum}
    E_{\mathrm{HB}}
    = \frac{D_0(\text{H}_2)}{\caKappaGeom^2 N_{\mathrm{shell}}(2)}
    = \frac{BE_H}{d\,\caKappaGeom^2 N_{\mathrm{shell}}(2)}
    = \frac{BE_H}{2^{\mathrm{frac\_bits}}-1}
    \approx 0{,}2160\ \mathrm{эВ}
    \approx 20{,}84\ \SiFrac{кДж}{моль}
  \end{equation}
  при длине моста донор--акцептор $R(\text{O}\cdots\text{O})\approx R(\text{N}\cdots\text{O}) = d\cdot R(\text{O--H})\approx 2{,}88\ \text{\normalfont\AA}$.
\item \textbf{Точная соизмеримость пептидного хребта с решёткой водорода:}
  по линейному закону ковалентных радиусов второго периода $r_{\mathrm{cov}}(Z)=\frac{39-2Z}{13}\caKappaGeom a_0$
  сумма ковалентных радиусов углерода ($Z=6$) и азота ($Z=7$) делится на кластерное число $N_\varphi=13$ \emph{нацело}:
  \begin{equation}
  \label{eq:bio-peptide-commensurability}
    R(\text{C--N})
    = r_{\mathrm{cov}}(6) + r_{\mathrm{cov}}(7)
    = \frac{27 + 25}{13}\,\caKappaGeom\,a_0
    = \frac{52}{13}\,\caKappaGeom\,a_0
    = 4\caKappaGeom\,a_0
    = 2 R_e(\text{H}_2)
    \approx 1{,}497\ \text{\normalfont\AA}.
  \end{equation}
\end{enumerate}

\section{Трёхдорожечный замок Уотсона--Крика в гексагональном срезе \texorpdfstring{$\{111\}$}{{111}}}
\label{sec:bio-watson-crick}

В плоском срезе $\{111\}$ решётки~\latGCKshort\ ($sp^2$, звезда~$N_6$ с углом~$120^\circ$) азотистые основания образуют жёсткие плоские гетероциклы: пурины ($\text{A, G}$: сшивка $6$- и $5$-членного колец) и пиримидины ($\text{C, T/U}$: $6$-членное кольцо).
Постоянный поперечник двойной спирали фиксируется суммой одного пурина и одного пиримидина, а вдоль фронтального стыка идут ровно $d=3$ параллельные дорожки водородных связей со знаками полярности ($+$ для донора $-\text{NH}$, $-$ для акцептора $=\text{O}$ или $=\text{N}-$, $0$ для неполярного $-\text{CH}$):
\begin{itemize}[noitemsep]
\item \textbf{Пара $\text{G--C}$:} $\text{G}=[-,\ +,\ +]$ $\longleftrightarrow$ $\text{C}=[+,\ -,\ -]$ --- замыкаются все $d=3$ дорожки с энергией $E(\text{G--C})=3\,E_{\mathrm{HB}}\approx 0{,}648\ \mathrm{эВ}$ ($14{,}94\ \SiFrac{ккал}{моль}$);
\item \textbf{Пара $\text{A--T}$ ($\text{A--U}$):} $\text{A}=[+,\ -,\ 0]$ $\longleftrightarrow$ $\text{T}=[-,\ +,\ -]$ --- замыкаются $d-1=2$ дорожки с энергией $E(\text{A--T})=2\,E_{\mathrm{HB}}\approx 0{,}432\ \mathrm{эВ}$ ($9{,}96\ \SiFrac{ккал}{моль}$), отношение $E(\text{G--C})/E(\text{A--T})=d/(d-1)=3/2$;
\item \textbf{Геометрический запрет некомплементарных пар ($\text{A--C}$, $\text{G--T}$):} столкновение одноимённых знаков ($++$ и $--$) на двух дорожках одновременно вызывает кулоновское отталкивание и выталкивает ошибочное основание из гексагонального среза $\{111\}$.
\end{itemize}

Вдоль оси двойной спирали B-ДНК вертикальный шаг $\pi$-стэкинга между соседними этажами пар оснований равен произведению битового бюджета ячейки $\lfloor\caBitBudget\rfloor=9$ на ковалентный радиус водорода $r_{\mathrm{cov}}(\text{H})=\caKappaGeom a_0$:
\begin{equation}
\label{eq:bio-dna-stacking-step}
  h_{\mathrm{stack}}(\text{B-DNA})
  = \lfloor\caBitBudget\rfloor\,\caKappaGeom\,a_0
  = 9\caKappaGeom\,a_0
  \approx 3{,}368\ \text{\normalfont\AA}
\end{equation}
(рентгеноструктурные данные B-ДНК: $3{,}38\text{--}3{,}40\ \text{\normalfont\AA}$, откл.\ $-0{,}36\%$).
Число пар оснований на один полный виток двойной спирали равно экваториальному коэффициенту квадратичного роста сфер кубооктаэдра $N_{\mathrm{shell}}(k)=(\geomNeighborTwelve-2)k^2+2$:
\[
  n_{\mathrm{turn}}(\text{B-DNA}) = \geomNeighborTwelve - 2 = 10\text{ пар оснований на виток}
\]
(шаг витка $P = 10\,h_{\mathrm{stack}}\approx 33{,}7\ \text{\normalfont\AA}$, поворот $360^\circ/10 = 36^\circ$),
а универсальный энергетический квант клетки --- свободная энергия гидролиза терминальной фосфатной связи АТФ ($\text{АТФ}\to\text{АДФ}+\text{Ф}_i$) --- равен удвоенному кванту водородной связи~\eqref{eq:bio-hbond-quantum} (равенству энергии двухдорожечной пары~$\text{A--T}$):
\[
  \Delta G_{\mathrm{ATP}}
  = (d-1)\,E_{\mathrm{HB}}
  = \frac{d-1}{d\,\caKappaGeom^2 N_{\mathrm{shell}}(2)}\,BE_H
  = \frac{d-1}{2^{\mathrm{frac\_bits}}-1}\,BE_H
  \approx 0{,}432\ \mathrm{эВ}
  \approx 41{,}7\ \SiFrac{кДж}{моль}
\]
(физиологический диапазон в клетке: $42\text{--}50\ \SiFrac{кДж}{моль}$).

\section{Четырёхслотовый штрихкод большой бороздки: чтение генома без расплетания спирали}
\label{sec:bio-major-groove-barcode}

В живой клетке транскрипционные факторы находят свои сайты связывания (промоторы, энхансеры, репрессоры) на интактной двойной спирали ДНК \emph{до} её расплетания РНК-полимеразой.
Это возможно потому, что тыльный край каждой пары Уотсона--Крика, обращённый в \textbf{большую бороздку (Major Groove)}, выставляет наружу уникальный четырёхпозиционный вектор из доноров ($+$), акцепторов ($-$), гидрофобного метильного шипа тимина ($M$) и нейтрального водорода ($0$):
\begin{equation}
\label{eq:bio-major-groove-code}
\begin{aligned}
  \text{A--T} &\longrightarrow [-,\ +,\ -,\ M], &\qquad
  \text{T--A} &\longrightarrow [M,\ -,\ +,\ -], \\
  \text{G--C} &\longrightarrow [-,\ -,\ +,\ 0], &\qquad
  \text{C--G} &\longrightarrow [0,\ +,\ -,\ -].
\end{aligned}
\end{equation}
Все четыре ориентированные пары ($\text{A--T}$, $\text{T--A}$, $\text{G--C}$, $\text{C--G}$) попарно различны в~\eqref{eq:bio-major-groove-code}.
Боковые радикалы аминокислот распознающей спирали белка несут комплементарные вилки квантов~$E_{\mathrm{HB}}$:
\begin{itemize}[noitemsep]
\item гуанидиновая вилка аргинина ($\text{Arg}: [+,\ +]$) бидентатно защёлкивается на двух соседних акцепторах гуанина ($[-,\ -]$ в $\text{G--C}$, выигрыш $2\,E_{\mathrm{HB}}$);
\item амидная вилка аспарагина и глутамина ($\text{Asn/Gln}: [+,\ -]\leftrightarrow [-,\ +]$) бидентатно защёлкивается на паре $[-,\ +]$ аденина ($\text{A--T}$, выигрыш $2\,E_{\mathrm{HB}}$);
\item карбоксилатная вилка глутамата и аспартата ($\text{Glu/Asp}: [-,\ -]$) фиксирует донор~$+$ цитозина ($\text{C--G}$) или аденина;
\item неполярный карман ($\text{Ala, Val, Leu, Ile, Thr}$) фиксирует метильный шип~$M$ тимина за счёт вытеснения напряжённой клетки воды.
\end{itemize}
Тем самым распознавание регуляторного участка ДНК белком сводится к локальной проверке комплементарности знаков в большой бороздке за $O(1)$ на шаг без глобального перебора.

\section{Два универсальных ДНК-связывающих замка факторов транскрипции: петля \texorpdfstring{$\text{C2H2}$}{C2H2} и штифт \texorpdfstring{$\text{Trp}$}{Trp} гомеодомена}
\label{sec:bio-tf-dbd-locks}

Прямой анализ белковых последовательностей факторов транскрипции в референсных геномах \emph{Drosophila melanogaster}, \emph{Caenorhabditis elegans} и \emph{Homo sapiens} показывает, что распознающие домены укладываются в два канонических геометрических замка кубооктаэдра:
\begin{enumerate}[leftmargin=*]
\item \textbf{Цинковый палец $\text{C2H2}$ ($\text{Cys}_1\text{--x}_{2\ldots 4}\text{--Cys}_2\text{--петля--His}_1\text{--x}_{3\ldots 5}\text{--His}_2$):}
  расстояние по цепи от второго цистеина~$\text{Cys}_2$ до первого гистидина~$\text{His}_1$ равно размеру одетого кластера $N_\varphi = 1 + \geomNeighborTwelve$, так что длина внутренней распознающей петли равна ровно числу вершин первой сферы кубооктаэдра:
  \[
    \text{His}_1 - \text{Cys}_2 = N_\varphi \quad\Longrightarrow\quad L_{\mathrm{ZF}}(\text{Cys}_2\ldots\text{His}_1) = \geomNeighborTwelve.
  \]
  Внутри петли из $\geomNeighborTwelve$ остатков три соседние пары оснований кодонного триплета ДНК ($d$ пар в большой бороздке) считываются вилками водородных связей на трёх позициях с шагом~$d$:
  \[
    \bigl\{\,|\square|,\quad \lfloor\caBitBudget\rfloor,\quad \geomNeighborTwelve\,\bigr\}.
  \]
\item \textbf{Гомеодомен и спираль «спираль--поворот--спираль» со штифтом $\text{Trp}$ ($\text{W}$):}
  распознающая $\alpha$-спираль, входящая в большую бороздку B-ДНК, центрирована на инвариантном штифте~$\text{Trp}$ ($\text{W}$), за которым стоят ароматический стабилизатор $\text{Phe/Tyr/Ile}$ ($\text{W}+1$), бидентатный замок $\text{Asn}$ ($[+,-]$ на смещении $\text{W}+d$, фиксирующий пару $\text{A--T}$ ядра промотора) и катионный замок $\text{Arg}$ ($[+,+]$ на смещении $\text{W}+d_{\mathrm{out}}(E)$).
  Индивидуальная специфичность фактора транскрипции задаётся главным переключателем на смещении $\text{W}+(d-1)$ ($\text{Lys}$ у передних морфогенов \texttt{bcd, oc, OTX2, CRX}; $\text{Ser}$ у \texttt{ey, PAX6}; $\text{Gln}$ у сегментных \texttt{Hox}-генов \texttt{Dfd, Scr, Antp, Ubx, abd-A, Abd-B, lin-39, mab-5, HOXA5..HOXD10}; $\text{Cys}$ у \texttt{unc-86, POU3F1/2}; $\text{His}$ у \texttt{CUX1/2}) совместно с боковыми пазами $\text{W}-1,\ \text{W}-d_{\mathrm{in}}(\square),\ \text{W}-|\square|$.
\end{enumerate}

\section{Почему основание --- жёсткая плитка, а хромосома квантуется катушками по \texorpdfstring{$N_{\mathrm{ball}}(d)=147$}{N\_ball(d)=147} пар оснований на \texorpdfstring{$\vevHierarchySteps=8$}{N\_hier=8} гистонах}
\label{sec:bio-nucleosome-147}

\begin{enumerate}[leftmargin=*]
\item \textbf{Почему буква генома не может быть одиночным атомом и не требует поатомного пересчёта:}
  при $300\ \mathrm{K}$ ($k_B T \approx 0{,}026\ \mathrm{эВ}$) одиночная водородная связь ($E_{\mathrm{HB}}\approx 8{,}3\,k_B T$) испытывала бы тепловой срыв с вероятностью $e^{-8{,}3}\approx 2{,}4\times 10^{-4}$ и не имела бы тыльной грани для чтения без разрыва цепи.
  Плоский гетероцикл азотистого основания в срезе $\{111\}$ скреплён внутренними ковалентными $sp^2$-связями ($\sim 4\text{--}7\ \mathrm{эВ}\gg k_B T$), которые при $300\ \mathrm{K}$ полностью заморожены.
  Поэтому на биологическом этаже $k=11$ каждое азотистое основание выступает как \textbf{одна монолитная жёсткая плитка}, разводящая на переднем торце $d=3$ водородные дорожки Уотсона--Крика, а на тыльном торце --- $|K_4|=4$ слота большой бороздки.
  Ширина пары оснований между гликозидными узлами $\text{C1}'\text{--C1}'$ равна
  \[
    R(\text{C1}'\text{--C1}')
    = \bigl(1 + d_{\mathrm{in}}(\square)\,d_{\mathrm{out}}(V)\bigr)\caKappaGeom\,a_0
    \approx 10{,}851\ \text{\normalfont\AA}
    \qquad\bigl(\text{рентген B-ДНК: }10{,}85\ \text{\normalfont\AA},\ \text{откл.\ }+0{,}01\%\bigr),
  \]
  а полный внешний диаметр двойной спирали B-ДНК равен
  \[
    D_{\mathrm{DNA}}
    = \bigl(N_{\mathrm{shell}}(2) + \geomNeighborTwelve\bigr)\caKappaGeom\,a_0
    \approx 20{,}21\ \text{\normalfont\AA}
    \qquad\bigl(\text{рентген B-ДНК: }20{,}0\ \text{\normalfont\AA},\ \text{откл.\ }+1{,}0\%\bigr).
  \]
\item \textbf{Квант нуклеосомной катушки хроматина ($N_{\mathrm{ball}}(d)=147$ пар оснований на $\vevHierarchySteps=8$ гистонах):}
  чтобы не сканировать миллиарды пар оснований на каждом такте, клетка наматывает неактивную нить ДНК на белковые катушки --- гистоновые октамеры из $\vevHierarchySteps = 2|K_4| = 8$ субъединиц ($2\times\{\text{H2A, H2B, H3, H4}\}$).
  Число пар оснований ДНК, намотанных на одну нуклеосомную катушку, в точности равно полному числу узлов трёхсферного шара кубооктаэдра~\latGCKshort\ ($k\le d$), в котором открываются все четыре канала $(V,E,\square,\triangle)$:
  \begin{equation}
  \label{eq:bio-nucleosome-147}
    N_{\mathrm{nuc}}
    = N_{\mathrm{ball}}(d)
    = \sum_{k=0}^{d} N_{\mathrm{shell}}(k)
    = 1 + \geomNeighborTwelve + N_{\mathrm{shell}}(2) + N_{\mathrm{shell}}(3)
    = 147\text{ пар оснований}
  \end{equation}
  (кристаллография нуклеосомы PDB \texttt{1AOI}: $147\ \mathrm{bp}$ на $8$ гистонах).
  Все катушки закрытого хроматина ($98\text{--}99\%$ длины хромосомы) экранируют большую бороздку и выключают её из тактового цикла, оставляя для $O(1)$-чтения только открытые регуляторные окна между катушками.
\end{enumerate}

\chapter{Геометрический вывод генетического кода: почему \texorpdfstring{$64\to 20+1=21$}{64 -> 20+1=21}}
\label{ch:bio-genetic-code}

\section{Откуда берутся \texorpdfstring{$64$}{64} кодона и \texorpdfstring{$21$}{21} смысловой класс}
\label{sec:bio-64-to-21}

В континуальной биологии числа универсального генетического кода --- $4$ буквы алфавита, длина кодона $3$, $64$ триплета и $20$ канонических аминокислот плюс $1$ стоп-сигнал ($21$ смысловой класс) --- принято объявлять «замороженной случайностью эволюции» (frozen accident).
На решётке~\latGCKshort\ все эти числа уже зафиксированы константами планковской ячейки и второй координационной сферы кубооктаэдра:
\begin{enumerate}[leftmargin=*]
\item \textbf{Четыре буквы алфавита ($|\mathcal{A}|=4=|K_4|$):}
  число вершин правильного тетраэдра~$K_4$ в~\latGCKshort\ и число различных штрихкодов~\eqref{eq:bio-major-groove-code} в большой бороздке равно~$4$ ($\{\text{U/T, C, A, G}\}$).
\item \textbf{Триплетность кодона ($d=3$) и $64=2^{\mathrm{frac\_bits}}$ состояния:}
  минимальное слово, фиксирующее точку в $d=3$ пространственных измерениях (три взаимно ортогональные плоскости клики~$K_3$), имеет длину $d=3$ буквы и порождает ровно
  \[
    4^d = 4^3 = 64 = 2^6 = 2^{\mathrm{frac\_bits}} = \frac{1}{z_{\min}}
  \]
  кодонов --- в точности число дискретных уровней амплитудной сетки планковской ячейки ($\mathrm{frac\_bits}=6$, $z_{\min}=1/64$).
\item \textbf{Емкость второй сферы кубооктаэдра ($N_{\mathrm{shell}}(2)/2 = 21 = 20 + 1$):}
  все ковалентные и водородные связи биомолекул живут на второй координационной сфере кубооктаэдра ($k=2$), содержащей
  \[
    N_{\mathrm{shell}}(2) = 10\cdot 2^2 + 2 = 42\text{ узла} = 12_V + 24_E + 6_\square.
  \]
  По принципу Паули (`A16`, 2 спиновых состояния на узел) вторая сфера кубооктаэдра вмещает ровно
  \begin{equation}
  \label{eq:bio-21-classes}
    \frac{N_{\mathrm{shell}}(2)}{2}
    = \frac{42}{2}
    = 21\text{ спиновую пару}
    = 6_V + 12_E + 3_\square,
  \end{equation}
  то есть ровно $20$ связующих классов (канонических аминокислот) и $1$ терминальный класс (\texttt{Stop})!
\end{enumerate}

\section{Симметрия Румера и тождество \texorpdfstring{$\vevHierarchySteps + N_\varphi = 8 + (12+1) = 21$}{N\_hier + N\_phi = 8 + (12+1) = 21}}
\label{sec:bio-rumer-symmetry}

Первые две позиции триплета ($X,Y\in\{\text{U, C, A, G}\}$) образуют $4^2=16=2\vevHierarchySteps$ дублетных боксов (`XY*`).
Согласно фундаментальной симметрии Румера, эти $16$ боксов делятся строго пополам на два октета по $\vevHierarchySteps=\lfloor\caBitBudget\rfloor-1=8$ боксов в каждом:
\begin{enumerate}[leftmargin=*]
\item \textbf{Октет~I ($\vevHierarchySteps=8$ нерасщеплённых семейных боксов):}
  в восьми боксах (\texttt{GC*}, \texttt{GG*}, \texttt{GU*}, \texttt{CC*}, \texttt{CU*}, \texttt{CG*}, \texttt{UC*}, \texttt{AC*}) третья буква кодона полностью вырождена (все $4$ кодона кодируют одну и ту же аминокислоту).
  Этот октет даёт ровно $\vevHierarchySteps = 8$ базовых аминокислот остова:
  \[
    \{\text{Ala, Gly, Val, Pro, Leu, Arg, Ser, Thr}\}\qquad (8\text{ аминокислот}).
  \]
\item \textbf{Октет~II ($\vevHierarchySteps=8$ расщеплённых боксов $\longrightarrow N_\varphi=1+\geomNeighborTwelve=13$ классов):}
  во втором октете третья позиция различает пурины ($\text{A, G}$) и пиримидины ($\text{U, C}$), порождая ровно $\geomNeighborTwelve=12$ специализированных аминокислот
  \[
    \{\text{Phe, Tyr, Cys, Trp, His, Gln, Asn, Lys, Asp, Glu, Ile, Met}\}\qquad (12\text{ аминокислот})
  \]
  плюс $1$ центральный класс обрыва цепи (\texttt{Stop}: \texttt{UAA, UAG, UGA}), то есть в точности $1+\geomNeighborTwelve = 13 = N_\varphi$ исходов одетого кластера!
\end{enumerate}
Складывая оба октета Румера, получаем фундаментальное структурное тождество генетического кода на решётке~\latGCKshort:
\begin{equation}
\label{eq:bio-rumer-identity}
  \underbrace{\vevHierarchySteps}_{\text{8 семейных АК}}
  + \underbrace{\bigl(\geomNeighborTwelve + 1\bigr)}_{N_\varphi\,=\,12\text{ расщеплённых АК}\,+\,1\text{ Stop}}
  = 8 + (12 + 1)
  = 20 + 1
  = 21
  = \frac{N_{\mathrm{shell}}(2)}{2}.
\end{equation}

При этом три осевых места канала квадратных граней $3_\square$ в разложении $21 = 6_V + 12_E + 3_\square$ отвечают ровно $d=3$ аминокислотам, которые служат мостом между Октетом~I и Октетом~II и имеют по $4+2=6$ кодонов каждая:
$\text{Leu}$ (якорный гидрофобный стержень ядра: \texttt{CU*} $+$ \texttt{UUR}),
$\text{Ser}$ (якорный полярный гидроксил: \texttt{UC*} $+$ \texttt{AGY})
и $\text{Arg}$ (главный бидентатный чтец большой бороздки ДНК: \texttt{CG*} $+$ \texttt{AGR}).
Полная классификация всех $21$ смысловых классов второй координационной сферы приведена в \cref{tab:bio-21-amino-acids}.

\begin{table}[htbp]
\centering
\caption{Все $21 = N_{\mathrm{shell}}(2)/2$ смысловых классов универсального генетического кода: разбиение Румера $8_{\text{Октет I}} + (12_{\text{Октет II}} + 1_{\text{Stop}}) = 21$, вилки водородных связей~$E_{\mathrm{HB}}$ и геометрическая роль в $O(1)$-фолдинге.}
\label{tab:bio-21-amino-acids}
\footnotesize
\setlength{\tabcolsep}{3.5pt}
\begin{tabular}{@{}llllll@{}}
\toprule
Октет & Класс & Кодоны & Паз $\text{C}_\alpha$ & Вилка $E_{\mathrm{HB}}$ & Роль на кубооктаэдре ГЦК \\
\midrule
I ($3_\square$) & Leu (L) & \texttt{CU*}, \texttt{UUR} (6) & $\alpha$ & $[0,0]$ (ветка $K_4$) & гидрофобный якорь ядра \\
I ($3_\square$) & Ser (S) & \texttt{UC*}, \texttt{AGY} (6) & $\gamma/\alpha$ & $[+,-]$ ($-\text{OH}$) & короткий полярный замок \\
I ($3_\square$) & Arg (R) & \texttt{CG*}, \texttt{AGR} (6) & $\alpha/\beta$ & $[+,+]$ (гуанидин) & бидентатный чтец $\text{G--C}$ бороздки \\
\midrule
I ($V/E$) & Ala (A) & \texttt{GC*} (4) & $\alpha$ & $[M]$ ($-\text{CH}_3$) & эталон правого паза $\alpha$-спирали \\
I ($V/E$) & Val (V) & \texttt{GU*} (4) & $\beta$ & $[M,M]$ ($C_\beta$-вилка) & эталон плоского паза $\beta$-листа \\
I ($V/E$) & Pro (P) & \texttt{CC*} (4) & $\gamma$ & цикл $\text{C}_\alpha\text{--N}$ & жёсткий излом петли (нет $-\text{NH}$) \\
I ($V/E$) & Gly (G) & \texttt{GG*} (4) & $\gamma$ & $[0]$ ($-\text{H}$) & гибкий шарнир без бокового радикала \\
I ($V/E$) & Thr (T) & \texttt{AC*} (4) & $\beta$ & $[+,-,M]$ & полярно-метильный паз $\beta$-листа \\
\midrule
II ($12_E/V$) & Phe (F) & \texttt{UUY} (2) & $\beta/\alpha$ & срез $\{111\}$ ($N_6$) & плоский ароматический стек \\
II ($12_E/V$) & Tyr (Y) & \texttt{UAY} (2) & $\beta$ & $\{111\} + [+,-]$ & ароматический стек с $-\text{OH}$ замком \\
II ($12_E/V$) & Trp (W) & \texttt{UGG} (1) & $\beta/\alpha$ & бицикл $\{111\} + [+]$ & крупный индольный клин ядра \\
II ($12_E/V$) & Cys (C) & \texttt{UGY} (2) & $\beta/\gamma$ & $-\text{SH}$ ($Z=16$) & ковалентная сшивка $\text{S--S}$ ($2{,}05\ \text{\normalfont\AA}$) \\
II ($12_E/V$) & Met (M) & \texttt{AUG} (1) & $\alpha$ & $-\text{S--CH}_3$ ($Z=16$) & старт-замок трансляции \\
II ($12_E/V$) & Ile (I) & \texttt{AUH} (3) & $\beta$ & $[M,M]$ (хиральный $C_\beta$) & плотная упаковка $\beta$-ядра \\
II ($12_E/V$) & His (H) & \texttt{CAY} (2) & $\alpha/\gamma$ & $[+ \leftrightarrow -]$ ($\mathrm{p}K_a\approx 6$) & переключатель протона при $300\ \mathrm{K}$ \\
II ($12_E/V$) & Gln (Q) & \texttt{CAR} (2) & $\alpha$ & $[+,-]$ (амид) & бидентатный чтец $\text{A--T}$ бороздки \\
II ($12_E/V$) & Asn (N) & \texttt{AAY} (2) & $\gamma/\alpha$ & $[+,-]$ (короткий амид) & чтец $\text{A--T}$ и шапка спирали \\
II ($12_E/V$) & Lys (K) & \texttt{AAR} (2) & $\alpha$ & $[+]$ ($-\text{NH}_3^+$) & солевой мост к фосфатам ДНК \\
II ($12_E/V$) & Asp (D) & \texttt{GAY} (2) & $\gamma/\alpha$ & $[-,-]$ ($-\text{COO}^-$) & фиксатор доноров и иона $\text{Mg}^{2+}$ \\
II ($12_E/V$) & Glu (E) & \texttt{GAR} (2) & $\alpha$ & $[-,-]$ (длинный $-\text{COO}^-$) & бидентатный акцептор и солевой мост \\
\midrule
II (ядро $1$) & \texttt{Stop} & \texttt{UAA,UAG,UGA} (3) & --- & освобождение цепи & обрыв трансляции и выход глобулы \\
\bottomrule
\end{tabular}
\end{table}

\chapter{Ко-трансляционный фолдинг белка за \texorpdfstring{$O(1)$}{O(1)} на шаг без парадокса Левинталя}
\label{ch:bio-o1-folding}

\section{Почему парадокс Левинталя --- артефакт континуума}
\label{sec:bio-levinthal-resolution}

Парадокс Левинталя возникает только тогда, когда углы вращения полипептидного остова $(\phi,\psi)$ считают непрерывными вещественными переменными, а радиусы атомов --- несоизмеримыми числами, требующими глобальной минимизации энергии по $\sim 3^L$ конформациям.
На решётке~\latGCKshort\ оба этих допущения ложны:
\begin{enumerate}[leftmargin=*]
\item \textbf{Дискретность узла $\text{C}_\alpha$ ($K_4$, $d=3$ паза):}
  пептидная группа ($\text{C}_\alpha\text{--C(=O)--NH--C}_\alpha$) жёстко фиксирована в плоском гексагональном срезе $\{111\}$ ($120^\circ$),
  а тетраэдрический узел~$\text{C}_\alpha$ ($sp^3$, $\cos\theta=-1/3$) имеет ровно $d=3$ дискретных ориентационных паза относительно предыдущего звена:
  \begin{itemize}[noitemsep]
  \item паз~$\alpha$ (спиральный поворот с замыканием водородного моста $i\to i-4$ на шаге $R(\text{O}\cdots\text{N})=d\cdot R(\text{O--H})\approx 2{,}88\ \text{\normalfont\AA}$);
  \item паз~$\beta$ (плоский зигзаг в срезе $\{111\}$ с латеральным замыканием водородных мостов между соседними нитями);
  \item паз~$\gamma/\text{turn}$ (разворот петли, принудительно выбираемый пролином $\text{Pro}$ с замкнутым циклом или глицином $\text{Gly}$).
  \end{itemize}
\item \textbf{Локальный выбор за $O(1)$ на выходе из рибосомного туннеля:}
  по мере синтеза цепи каждый новый остаток~$i$ взаимодействует только со своей локальной окрестностью кубооктаэдра ($\geomNeighborTwelve$ направлений) и выбирает тот из $d=3$ пазов, который даёт максимальное число замкнутых квантов $E_{\mathrm{HB}}=BE_H/(d\,\caKappaGeom^2 N_{\mathrm{shell}}(2))$ и вытесняет гидрофобные радикалы внутрь ядра от сетки воды.
  Число проверок на один остаток ограничено константой $O(1)$ ($d\cdot\geomNeighborTwelve$ элементарных операций), а полная укладка белка длины~$L$ выполняется за $O(L)$ тактов в реальном времени без глобального перебора.
\end{enumerate}

\section{Геометрия вторичных структур Полинга--Кори (\texorpdfstring{$\alpha$}{alpha}-спираль \texorpdfstring{$3{,}6_{13}$}{3.6\_13} и \texorpdfstring{$\beta$}{beta}-лист) из констант кубооктаэдра}
\label{sec:bio-pauling-corey}

Все геометрические параметры классических вторичных структур белка Полинга--Кори вытекают напрямую из параметров каналов кубооктаэдра~\latGCKshort\ без подгоночных коэффициентов:
\begin{enumerate}[leftmargin=*]
\item \textbf{Правая $\alpha$-спираль ($3{,}6_{13}$-спираль):}
  карбонильный кислород $\text{C}'_i=\text{O}_i$ остатка~$i$ замыкает водородный мост $R(\text{O}\cdots\text{N})=d\cdot R(\text{O--H})$ с амидным азотом $\text{N}_{i+d_{\mathrm{in}}(\square)}\text{--H}$ через шаг $d_{\mathrm{in}}(\square)=4$ остатка ($i\to i+4$), образуя замкнутое кольцо ровно из $N_\varphi = 1 + \geomNeighborTwelve = 13$ атомов.
  Полный период из $d_{\mathrm{out}}(E)=5$ витков спирали содержит $\geomNeighborTwelve + |\square| = 12 + 6 = 18$ остатков, откуда число остатков на один виток равно
  \[
    n_\alpha = \frac{\geomNeighborTwelve + |\square|}{d_{\mathrm{out}}(E)} = \frac{18}{5} = 3{,}6\text{ остатка/виток}
    \qquad\bigl(\text{угол поворота }2\pi/n_\alpha = 100^\circ\bigr).
  \]
  Осевой подъём на один остаток равен длине связи $R(\text{C--N}) = d_{\mathrm{in}}(\square)\,\caKappaGeom\,a_0$:
  \[
    h_\alpha = d_{\mathrm{in}}(\square)\,\caKappaGeom\,a_0 \approx 1{,}497\ \text{\normalfont\AA}
    \qquad\bigl(\text{рентген Полинга--Кори: }1{,}500\ \text{\normalfont\AA},\ \text{откл.\ }-0{,}22\%\bigr),
  \]
  шаг одного витка равен
  \[
    P_\alpha = n_\alpha\,h_\alpha = \frac{d\,\delta_{\mathrm{pack}}^{-1}}{d_{\mathrm{out}}(E)}\,\caKappaGeom\,a_0 \approx 5{,}388\ \text{\normalfont\AA}
    \qquad\bigl(\text{рентген: }5{,}400\ \text{\normalfont\AA},\ \text{откл.\ }-0{,}22\%\bigr),
  \]
  а длина полного $d_{\mathrm{out}}(E)$-виткового ($\geomNeighborTwelve+|\square|=18$ остатков) трансмембранного стержня совпадает с толщиной гидрофобного ядра липидного бислоя клеточной мембраны:
  \[
    L_{\mathrm{TM}} = \bigl(\geomNeighborTwelve + |\square|\bigr)\,h_\alpha = d\,\delta_{\mathrm{pack}}^{-1}\,\caKappaGeom\,a_0 \approx 26{,}94\ \text{\normalfont\AA}
  \]
  (гидрофобное ядро бислоя: $27{,}0\ \text{\normalfont\AA}$, откл.\ $-0{,}22\%$).
\item \textbf{Складчатый $\beta$-лист в срезе $\{111\}$:}
  осевой шаг на один остаток вдоль антипараллельной $\beta$-нити равен
  \[
    h_\beta = \left(\lfloor\caBitBudget\rfloor + \frac{1}{d}\right)\caKappaGeom\,a_0 \approx 3{,}492\ \text{\normalfont\AA}
    \qquad\bigl(\text{рентген Полинга--Кори: }3{,}48\ \text{\normalfont\AA},\ \text{откл.\ }+0{,}36\%\bigr).
  \]
\end{enumerate}

\section{Синаптический иммуноглобулиновый замок (\texorpdfstring{$\text{Ig}_{\mathrm{donor}}\leftrightarrow\text{Ig}_{\mathrm{acceptor}}$}{Ig\_donor <-> Ig\_acceptor})}
\label{sec:bio-ig-synaptic-lock}

В рецепторах клеточной адгезии, формирующих синаптическую специфичность нейронов (\texttt{Dpr} $\leftrightarrow$ \texttt{DIP}), длина замкнутой дисульфидной петли первого иммуноглобулинового домена $[\text{Cys}_B\ldots\text{Cys}_F]$ у донорных молекул ($\text{Ig}_{\mathrm{donor}}$) равна водородному делителю кубооктаэдра:
\[
  L_{\mathrm{loop}}(\text{Ig}_{\mathrm{donor}})
  = d\,\caKappaGeom^2 N_{\mathrm{shell}}(2)
  = \frac{1-\caKappaGeom^6}{\caKappaGeom^6}\,\lfloor\caBitBudget\rfloor
  = 2^{\mathrm{frac\_bits}} - 1,
\]
а междоменный шаг от $\text{Cys}_F$ первой петли до цистеина второго домена у донорного выступа (с гребнем $\text{Trp--x--Arg}$, вилка $[+,+]$) и у акцепторной впадины (с канавкой $\text{Trp--x--x--x--[Asp/Glu]}$, вилка $[-,-]$) отклоняется от числа узлов второй координационной сферы $N_{\mathrm{shell}}(2)$ ровно на~$\pm 1$:
\[
  \Delta_{\mathrm{Cys}}(\text{Ig}_{\mathrm{donor}}) = N_{\mathrm{shell}}(2) - 1 = \Delta\varphi_{\mathrm{disc}},\qquad
  \Delta_{\mathrm{Cys}}(\text{Ig}_{\mathrm{acceptor}}) = N_{\mathrm{shell}}(2) + 1,
\]
\[
  \Delta_{\mathrm{Cys}}(\text{Ig}_{\mathrm{donor}}) + \Delta_{\mathrm{Cys}}(\text{Ig}_{\mathrm{acceptor}}) = 2\,N_{\mathrm{shell}}(2).
\]
При антипараллельной стыковке $\beta$-граней Ключа и Замка в гексагональном срезе $\{111\}$ избирательность синаптических пар ($\Delta G_{\mathrm{bind}} = -\Delta n_{\mathrm{HB}}\,E_{\mathrm{HB}}$) определяется без внешних таблиц сродства комплементарностью вилок водородных связей ($[+]\leftrightarrow [-]$) и стерических карманов на семи структурных смещениях кубооктаэдра от консервативного штифта $\text{Trp}$ ($\text{W}$) и конца петли $\text{Cys}_F$:
\[
  \bigl\{\text{W}+d,\ \ \text{W}+d_{\mathrm{in}}(\square),\ \ \text{W}+|\square|,\ \ \text{W}+\geomNeighborTwelve,\ \ \text{W}+N_\varphi,\ \ \text{Cys}_F-\geomNeighborTwelve,\ \ \text{Cys}_F-d_{\mathrm{out}}(E)\bigr\}:
\]
\begin{itemize}[noitemsep]
\item \textbf{Контакт $\text{W}+d_{\mathrm{in}}(\square)$ и $\text{W}+|\square|\leftrightarrow \text{W}+d$:}
  вилка $\text{Arg}$ ($[+,+]$ на $\text{W}+d_{\mathrm{in}}(\square)$) и гидрофобный $\text{Leu}$ ($[M,M]$ на $\text{W}+|\square|$) связывают $\text{Asp}$ ($[-,-]$ на $\text{W}+d_{\mathrm{in}}(\square)$) и $\text{Val}$ ($\text{W}+d$) группы \texttt{DIP-theta/eta}, тогда как пара $\text{Lys}$ ($[+]$ на $\text{W}+d_{\mathrm{in}}(\square)$) $+$ ароматическое кольцо $\text{Phe/Tyr}$ ($\{111\}$ на $\text{W}+|\square|$) избирательно защёлкивается на $\text{Glu}$ ($[-,-]$ на $\text{W}+d_{\mathrm{in}}(\square)$) $+$ $\text{Phe}$ ($\text{W}+d$) группы \texttt{DIP-epsilon/zeta};
\item \textbf{Контакт $\text{W}+\geomNeighborTwelve$:}
  отрицательный акцептор $\text{Asp/Glu}$ ($[-,-]$ на $\text{W}+\geomNeighborTwelve$ у \texttt{dpr11/15/16}) защёлкивается на положительной гребёнке $\text{Arg/Lys}$ ($[+,+]$ на $\text{W}+N_\varphi+1$ и грани~$F$) замка \texttt{DIP-gamma};
\item \textbf{Контакт $\text{Cys}_F-\geomNeighborTwelve \leftrightarrow \text{Cys}_F-d_{\mathrm{out}}(E)$ и $\text{W}+N_\varphi$:}
  индольный клин $\text{Trp}$ на $\text{Cys}_F-\geomNeighborTwelve$ и гидроксил $\text{Ser/Thr}$ ($[+,-]$ на $\text{W}+N_\varphi$ у \texttt{dpr6/10}) точно входят в полость $\text{Gly}$ на $\text{Cys}_F-d_{\mathrm{out}}(E)$ и акцептор $\text{Glu}$ ($\text{W}+N_\varphi$) замка \texttt{DIP-alpha}, тогда как $\text{Tyr}$ на $\text{Cys}_F-\geomNeighborTwelve$ и $\text{Arg}$ на $\text{W}+N_\varphi$ (\texttt{dpr8/9}) фиксируют замок \texttt{DIP-beta}.
\end{itemize}

\chapter{От молекулярных замков к клеточному автомату и 3D-онтогенезу}
\label{ch:bio-ontogeny-ca}

\section{Пошаговый цикл клетки без захардкоженных тегов}
\label{sec:bio-cell-cycle-o1}

Построенная лестница ($\text{этажи }0\text{--}2 \to \text{химия }Z=1\ldots 16 \to E_{\mathrm{HB}}=BE_H/(d\,\caKappaGeom^2 N_{\mathrm{shell}}(2)) \to \text{штрихкод большой бороздки} \to O(1)\text{-фолдинг}$) задаёт прямую архитектуру клеточного автомата морфогенеза (\texttt{genome-printer} / \texttt{CASA}) на макро-решётке кубооктаэдра~$N_{12}$:
\begin{enumerate}[leftmargin=*]
\item \textbf{Трансляция и $O(1)$-укладка:}
  кодирующая последовательность гена переводится по таблице~\eqref{eq:bio-rumer-identity} ($64\to 20+1$) в полипептидную цепь, которая ко-трансляционно за $O(L)$ шагов укладывается на решётке кубооктаэдра, формируя на поверхности глобулы распознающую вилку доноров и акцепторов~$E_{\mathrm{HB}}$.
\item \textbf{Цис-регуляторное чтение без строковых подгонок:}
  свёрнутый белок-регулятор сканирует четырёхслотовый штрихкод большой бороздки~\eqref{eq:bio-major-groove-code} в регуляторном окне гена; связывание и запуск транскрипции определяются суммой замкнутых водородных квантов $\sum n_{\mathrm{HB}}\,E_{\mathrm{HB}}$ выше теплового порога $8\,k_B T_{300\mathrm{K}}$, а не внешними текстовыми метками \texttt{GFF3}.
\item \textbf{Локальный шаг ткани на кубооктаэдре $N_{12}$:}
  каждая клетка ткани занимает макро-узел ромбического додекаэдра с $12$ соседями~$N_{12}$, обменивается морфогенами по тем же $12$ каналам с весом $\kappa_{\mathrm{link}}=1/12$ и принимает решение о делении, апоптозе или дифференцировке строго локально за $O(1)$ операций на такт.
\end{enumerate}

\input{volumes/include-computational-laboratory.tex}

\input{volumes/include-measurement-processing.tex}
